% Independent mathematical derivation. Prefix FLW.
% Source attribution and scope are part of the proof, not optional metadata.
\section{The full-lattice projective Verschiebung and its retained cycle object}
\label{sec:full-lattice-witt-receiver}

Alain Connes and Caterina Consani define the Verschiebung of an
endomorphism by cyclically moving its summands, with the original
endomorphism at the returning corner, in
\href{https://arxiv.org/abs/2004.08879v1}{\emph{Segal's Gamma rings
and universal arithmetic}}, original author source
\nolinkurl{Froblifts.tex}, lines 454--681, especially
\texttt{defninv}, \texttt{w0s}, \texttt{rabi9}, \texttt{rabi10},
\texttt{rhoinf}, and \texttt{lognewton}.
Their construction uses finite pointed sets and their associated
$\mathbb S$-modules. The construction below is instead carried out
on explicit finitely generated projective semimodules over the
programme's full $G_L(\mathbb C)$. We prove its coefficient maps,
the intrinsic trace, and the actual Frobenius/Verschiebung relations;
we then construct an exact receiving map to the authors' pointed
cycle object. Neither scalar group completion nor a single lattice
trace is assumed to retain all projective support.
The cited source expressly credits Gert Almkvist for the original
endomorphism invariant and Michael Atiyah for the $\lambda$-ring
framework used with Adams operations. Those earlier papers are
human precedents cited through this source; no fresh reading of
their complete proofs is claimed here.

The original complex companion, its full metric, and its signed
odd receiver are TPT1--TPT28 and TOB6--TOB29 of the present reader.
The full carrier and its coordinate-mask correspondence are
OLR16--OLR17. Their predecessor source is
\href{https://github.com/KokunoYumeto/zeta-function-research-reader/blob/3c78cfbd9e194d35117b33277c421c6564ce3dbf/workbenches/splitzero-tandem/continuations/20260923-full-prime-support/033/MIXED_PRIME_ENERGY_AND_BOUNDARY.tex}{033,
SBT1--SBT23}. The earlier supported-zero prime construction retains
its original provenance; it is not a discovery of this section.

\subsection{The full coefficient semiring and projective carriers}

Let $L$ be a bounded distributive lattice, with $0_L\ne1_L$, and put
\[
\begin{gathered}
 S=G_L(\mathbb C)
   =\{(c,1_L):c\in\mathbb C\}\cup\{z_\lambda=(0,\lambda):\lambda\in L\},\\
 (a,\lambda)+(b,\mu)=(a+b,\lambda\vee\mu),\\
 (a,\lambda)(b,\mu)=(ab,\lambda\wedge\mu),\qquad
 \tau=z_{0_L},\quad e=z_{1_L},\quad1_S=(1,1_L).
\end{gathered}
\tag{FLW1}
\]
The two coefficient maps
\[
 \rho:S\longrightarrow\mathbb C,\quad(c,\lambda)\longmapsto c,
 \qquad
 s:S\longrightarrow L,\quad(c,\lambda)\longmapsto\lambda
 \tag{FLW2}
\]
are unital semiring homomorphisms, by their displayed component
laws. They are jointly injective. Their joint image consists exactly
of the pairs with $c\ne0\Rightarrow\lambda=1_L$.
In particular $\rho(\tau)=\rho(e)=0$, whereas
$s(\tau)=0_L$ and $s(e)=1_L$.

For a positive integer $r$, let $J_r$ be the all-$1_L$ support
matrix and let
\[
 \mathsf P_r=(I_r,J_r)\in M_r(S),\qquad
 P_r=\operatorname{im}\mathsf P_r
 =\{(v,\lambda):v\in\mathbb C^r,\ 
                       v\ne0\Rightarrow\lambda=1_L\}.
 \tag{FLW3}
\]
Here a matrix pair records its amplitude and support, with every
nonzero amplitude entry forcing support $1_L$; matrix multiplication
uses ordinary complex multiplication in amplitude and join-of-meets
in support. On a free vector $(v,\mu_1,\ldots,\mu_r)$,
$\mathsf P_r$ leaves $v$ unchanged and sets every support coordinate
to $\bigvee_j\mu_j$. Thus $\mathsf P_r^2=\mathsf P_r$ and its
image is exactly FLW3. This proves that $P_r$ is projective, with
the indicated splitting of the free module $S^r$.
For any complex $b$-by-$a$ matrix $T$, the matrix
$(T,J_{b,a})$ on the retracts gives
\[
 T^\uparrow:P_a\longrightarrow P_b,
       \qquad(v,\lambda)\longmapsto(Tv,\lambda).
 \tag{FLW4}
\]
It is $S$-linear by FLW1, including cancellation of amplitudes.
Composition and addition preserve this lift, and an invertible
square $T$ gives an automorphism with inverse $(T^{-1})^\uparrow$.
The lift of the complex zero map is $(v,\lambda)\mapsto(0,\lambda)$,
which is not the constant-$\tau$ map.

The source dimension throughout is $q=(k+1)^2>0$. The actual cyclic
carrier with $d$ slots is the direct sum of $d$ such projectives:
\[
 P_{q,d}=P_q^{\oplus d}=\operatorname{im}
        \operatorname{diag}(\mathsf P_q,\ldots,\mathsf P_q).
 \tag{FLW5}
\]
An element has independent slot supports
$((v_0,\lambda_0),\ldots,(v_{d-1},\lambda_{d-1}))$, with
$v_j\ne0\Rightarrow\lambda_j=1_L$.
For the original free coordinate masks the complete fibre is
\[
 \{(\mu_{j,i}):\bigvee_{i=1}^q\mu_{j,i}=\lambda_j
                 \text{ for each }j,
             \ (v_j)_i\ne0\Rightarrow\mu_{j,i}=1_L\}.
 \tag{FLW6}
\]
This follows by applying FLW3 separately to every slot. The tensor
of two such synchronized projectives is also explicit:
$P_a\otimes_S P_b=\operatorname{im}(\mathsf P_a\otimes\mathsf P_b)$
under the standard free-coordinate identification, and the matrix
there is $(I_{ab},J_{ab})$. Hence it is $P_{ab}$.
This assertion follows directly by tensoring the two splitting maps:
their retraction composite is the identity and their opposite
composite is the displayed Kronecker idempotent.

\subsection{Intrinsic trace, including both zeros}

For any finitely generated projective $S$-semimodule $P$, choose
a splitting $i:P\to S^a$, $r:S^a\to P$, $ri=I_P$.
For $T:P\to P$ define
\[
 \operatorname{Tr}_P(T)=\operatorname{tr}_{S^a}(iTr)
                     =\sum_{j=1}^a(iTr)_{jj}\in S.
 \tag{FLW7}
\]
This is independent of the chosen splitting. Indeed for a second
splitting $i',r'$ through $S^b$, put $U=iTr'$ and $V=i'r$.
Then $UV=iTr$ and $VU=i'Tr'$. For rectangular matrices over a
commutative semiring,
$\operatorname{tr}(UV)=\sum_{j,l}U_{jl}V_{lj}
=\sum_{l,j}V_{lj}U_{jl}=\operatorname{tr}(VU)$.
This proves independence without subtraction or a chosen basis of
the projective. The same proof gives cyclicity for maps between two
projectives, invariance under isomorphism, additivity on direct sums,
and multiplicativity on tensor products, the last by the diagonal
entries of a Kronecker product.

For the chosen projective lift FLW4, every diagonal support is top,
so its intrinsic trace is exactly
\[
 \operatorname{Tr}_{P_q}(T^\uparrow)
       =(\operatorname{tr}_{\mathbb C}T,1_L).
 \tag{FLW8}
\]
In particular the supported zero map has trace $e$ and the actual
zero endomorphism has trace $\tau$. These are distinct even when the
complex trace is zero. More generally, applying $\rho$ or $s$ to
the splitting matrices gives the corresponding split projective
over $\mathbb C$ or $L$, and
\[
 \rho(\operatorname{Tr}_P T)
       =\operatorname{Tr}_{P\otimes_S\mathbb C}T_\rho,
 \qquad
 s(\operatorname{Tr}_P T)
       =\operatorname{Tr}_{P\otimes_S L}T_s.
 \tag{FLW9}
\]
One may define each base change here as the image of the coefficient
image of its splitting idempotent; the split maps prove the stated
tensor descriptions. FLW9 then follows by applying the coefficient
homomorphism to each finite diagonal sum. In particular
$P_q\otimes_S\mathbb C=\mathbb C^q$ and
$P_q\otimes_S L=\operatorname{im}J_q\cong L$, the latter by
evaluation at any coordinate of the constant support vector.

\subsection{The actual signed cyclic operator and every integer power}

Keep every original parameter and metric:
\[
\begin{gathered}
 k\ge17,\quad k\equiv1\pmod4,\quad\delta,\gamma>0,
 \quad q=(k+1)^2,\quad q-1\le N\le2q,\\
 E=\mathbb C[y]/\prod_{a,b=0}^k
      (y-(2b-k)\gamma+i(2a-k)\delta),\quad M[v]=[yv],\\
 F_p=\exp(i(\log p)M),\quad\widetilde F_p=p^{k/2}F_p,
 \quad\langle v,w\rangle_N=v^*G_Nw,\\
 \omega_{ab}=p^{(2a-k)\delta}
                 \exp(i(2b-k)\gamma\log p).
\end{gathered}
\tag{FLW10}
\]
For every positive odd $d$ retain
$\chi=\chi_d=(-1)^{(d-1)/2}$, $A=\chi F_p^\chi$.
On the carrier FLW5 define the literal Verschiebung
\[
 \mathcal C_d(A)((v_j,\lambda_j)_{j=0}^{d-1})
 =((Av_{d-1},\lambda_{d-1}),
                       (v_0,\lambda_0),\ldots,(v_{d-2},\lambda_{d-2})).
 \tag{FLW11}
\]
The ambient block matrix has $A^\uparrow$ in its corner,
$I^\uparrow$ on its cyclic forward edges, and the actual
all-$\tau$ zero matrix in all other blocks. Thus its support is
$P_d\otimes J_q$, where $P_d$ is the cyclic permutation support
matrix; it is not the all-top $dq$-by-$dq$ matrix.
Its inverse reverses the cyclic steps and uses $(A^{-1})^\uparrow$
at the reversed corner. Consequently it is an automorphism of
$P_{q,d}$, whose identity is the block idempotent in FLW5.

For every integer $n$ and every input slot $j\in\{0,\ldots,d-1\}$,
write $j+n=dw+r$, $0\le r<d$, using Euclidean division even for
negative $n$. The power map has exactly the following one block
in that input column:
\[
 \mathcal C_d(A)^n\iota_j=\iota_r(A^w)^\uparrow,
 \qquad r=j+n\bmod d,\quad w=\left\lfloor\frac{j+n}{d}\right\rfloor.
 \tag{FLW12}
\]
Following one forward step increments the slot and adds one to the
exponent only at a wrap. This proves FLW12 for $n\ge0$ by induction.
The stated inverse proves the backward step and hence all $n<0$.
Its support matrix is $P_d^n\otimes J_q$ at every integer $n$,
including $n=0$, which gives the projective identity in FLW5.
The exact intrinsic power traces are therefore
\[
 \boxed{\operatorname{Tr}_{P_{q,d}}\mathcal C_d(A)^n
 =\begin{cases}
 \tau,&d\nmid n,\\
 \left(d\chi^m\displaystyle\sum_{a,b=0}^k
                   \omega_{ab}^{\chi m},1_L\right),&n=dm.
 \end{cases}}
 \tag{FLW13}
\]
At a nonreturn time, every diagonal block is the actual zero block,
so the result is $\tau$. At a return time $n=dm$, every diagonal
block is $(A^m)^\uparrow$; FLW8 and the ordinary eigenvalue trace
give the second line. For $n=0$ it is $(dq,1_L)$.
For cancellation of its complex power sum it is $e$, not $\tau$.
No metric is involved in this intrinsic trace; in particular none
of the nonorthogonal Lagrange vectors is silently made orthogonal.

Let $P_{ab}$ be the actual original Lagrange projector and lift
its block diagonal as $\mathcal P_{ab}$ on $P_{q,d}$. The labelled
and mass-weighted identities are
\[
\begin{split}
 \operatorname{Tr}(\mathcal P_{ab}\mathcal C_d(A)^n)
 &=\begin{cases}\tau,&d\nmid n,\\
       (d\chi^m\omega_{ab}^{\chi m},1_L),&n=dm,
    \end{cases}\\
 \operatorname{Tr}(\mathcal P_{ab}(p^{k/2},1_L)^n\mathcal C_d(A)^n)
 &=(p^{kn/2},1_L)
              \operatorname{Tr}(\mathcal P_{ab}\mathcal C_d(A)^n).
\end{split}
\tag{FLW14}
\]
These follow from $P_{ab}A^m=\chi^m\omega_{ab}^{\chi m}P_{ab}$
and $\operatorname{tr}P_{ab}=1$. The projectors are retained as
specific maps: for distinct labels their supported product is the
lift of the complex zero map on each slot, not the actual zero
endomorphism. Thus their complex direct-sum decomposition has not
been incorrectly declared an orthogonal projective decomposition
over $S$.

The two full coefficient images of FLW13, as formal ghost series
at positive times, are
\[
\begin{split}
 \mathfrak g_d(t)&=\sum_{n\ge1}
                  \operatorname{Tr}(\mathcal C_d(A)^n)t^n\in S[[t]],\\
 \rho(\mathfrak g_d(t))
     &=d\sum_{m\ge1}\chi^m
                       \Big(\sum_{a,b}\omega_{ab}^{\chi m}\Big)t^{dm},\\
 s(\mathfrak g_d(t))&=\sum_{m\ge1}1_L t^{dm}\in L[[t]].
\end{split}
\tag{FLW15}
\]
All absent coefficients of the first and last series are their
respective additive zeros. The mass-weighted series inserts the
literal $p^{kdm/2}$ into the second line and leaves the third line
unchanged. These are formal series, so no analytic convergence is
being substituted for a finite coefficient calculation. The support
trace records divisibility of the return time but not its integer
multiplicity $d$; addition in $L$ is idempotent. The stronger
cyclic-support object is retained below.

\subsection{The unchanged metric, including support left by cancellation}

The full block pairing of two vectors in FLW5 is
\[
 \mathfrak h(v,w)=
   \left(\sum_{j=0}^{d-1}v_j^*G_Nw_j,
                     \bigvee_{j=0}^{d-1}(\lambda_j\wedge\mu_j)\right).
 \tag{FLW16}
\]
It takes values in $S$: a nonzero summand in its complex coordinate
requires both involved supports to be top. It is the sum of the
original supported Hermitian pairings, with the involution conjugating
amplitude and fixing support. Under FLW11 the supports are merely
permuted, so the exact signed difference formed using $(-1,1_L)$ is
\[
\begin{split}
 &\mathfrak h(\mathcal C_d(A)v,\mathcal C_d(A)w)
                  +(-1,1_L)\mathfrak h(v,w)\\
 &\qquad=
 \left(v_{d-1}^*(A^*G_NA-G_N)w_{d-1},
                     \bigvee_j(\lambda_j\wedge\mu_j)\right).
\end{split}
\tag{FLW17}
\]
The complex cancellation is the original TOB13 identity. Its support
does not cancel: the zero-amplitude contributions from the other
slots still supply the displayed join. Replacing that join by only
$\lambda_{d-1}\wedge\mu_{d-1}$ would be a false semiring equality.
For the mass-weighted operator the complex defect against $p^k$
times the original pairing is $p^k$ times the same displayed
amplitude, while its support remains this same join. All phases
and all entries of $G_N$ in $A^*G_NA$ are unchanged.

\subsection{Exact Frobenius and Verschiebung maps on the projectives}

For an automorphism $T$ of a projective $P$ define
$F_n(P,T)=(P,T^n)$ for $n\ge1$, and let $V_d(P,T)$ be its cyclic
construction on $P^{\oplus d}$ with corner $T$. We use actual zero
off-block maps in this definition. These definitions apply in
particular to FLW11 with $P=P_q$ and $T=A^\uparrow$.
The equalities $F_mF_n=F_{mn}$ follow by powers.
The isomorphism $V_mV_n(P,T)\cong V_{mn}(P,T)$ sends the outer
slot $i$ and inner slot $j$ to $i+mj$. One outer wrap advances the
inner slot and the final wrap uses $T$, so the permutation proves
the identity on every coordinate. It preserves the unchanged direct
sum metric and all independent slot supports.

For the full mixed relation let $g=\gcd(d,n)$, $l=d/g$, and
$a=n/g$. There is an explicit isomorphism
\[
 F_nV_d(P,T)\cong
        \bigoplus_{t=0}^{g-1}V_l(P,T^a).
 \tag{FLW18}
\]
For $P=P_q$ and $T=A^\uparrow$, its map from the right to the left
sends the slot $j$ of the $t$th summand to the original slot
$r_{t,j}=t+jn\bmod d$, with block map $(A^{w_{t,j}})^\uparrow$,
where $w_{t,j}=\lfloor(t+jn)/d\rfloor$ and $0\le j<l$.
The pairs $(t,j)$ enumerate every slot exactly once: the congruence
classes modulo $g$ are the disjoint orbits, each of length $l$.
The coordinate action FLW12 adds exponent
$w_{t,j+1}-w_{t,j}$ at each step. All intermediate steps therefore
intertwine, and $w_{t,l}=a$, $r_{t,l}=t$ proves the returning corner
$A^a$. Invertibility of $A$ proves invertibility of this map.
For a general invertible $T$ use its powers in the same formula.
In particular $F_dV_d(P,T)=(P,T)^{\oplus d}$ and, when $(d,n)=1$,
$F_nV_d(P,T)\cong V_dF_n(P,T)$.

The amplitude comparison is not asserted to be an isometry for
$G_N$ when powers of $A$ occur. Its exact pulled-back form is
\[
 \sum_{t=0}^{g-1}\sum_{j=0}^{l-1}
   x_{t,j}^*(A^{w_{t,j}})^*G_N A^{w_{t,j}}y_{t,j}.
 \tag{FLW19}
\]
Its support pairing is the unchanged corresponding join, because
each block map preserves its support. Thus the algebraic relation
has its complete metric comparison and does not remove a weight
defect by changing metrics.

For automorphisms $X$ of $P$ and $Y$ of $Q$, with product defined by
tensoring, the projection formula is the concrete isomorphism
\[
 V_n(P\otimes Q,X^n\otimes Y)
                 \cong(P,X)\otimes V_n(Q,Y).
 \tag{FLW20}
\]
From left to right use $X^j\otimes I$ in slot $j$. On the right
each nonreturning step uses $X\otimes I$ and the return uses
$X\otimes Y$; at the last step their product is $X^n\otimes Y$.
This checks the intertwining and gives its inverse $X^{-j}\otimes I$.
Likewise
\[
 V_n(P,X)\otimes V_n(Q,Y)
                 \cong\bigoplus_{t=0}^{n-1}V_n(P\otimes Q,X\otimes Y).
 \tag{FLW21}
\]
On the $t$th summand send slot $j$ to the pair
$(j,t+j\bmod n)$, using block map
$I\otimes Y^{\lfloor(t+j)/n\rfloor}$.
The simultaneous cycle steps preserve $t$; when either coordinate
wraps, its $X$ or $Y$ is exactly supplied by this formula, and the
full return uses $X\otimes Y$. The orbits partition all $n^2$
slots and all block maps are invertible, proving FLW21.
For $P=P_a$, $Q=P_b$ with their retained complex metrics
$G_P,G_Q$, the exact metric comparisons are respectively
$(X^j)^*G_PX^j\otimes G_Q$ and
$G_P\otimes (Y^{\lfloor(t+j)/n\rfloor})^*G_Q
                         Y^{\lfloor(t+j)/n\rfloor}$.
The support maps are permutations of the actual slot labels,
with the indicated tensor meet within each slot. No cancellation
of support or unproved unitary conjugacy is used.

The trace ghosts obey these isomorphisms in $S$, by FLW7. For example
FLW18 says that the $r$th trace of its left side is the sum of $g$
copies of the $r$th trace of $V_l(P,T^a)$.
For FLW11 the support condition on both sides is
$d\mid nr$, equivalently $l\mid r$; at return the complex
multiplicity is $gl=d$. Thus the integer multiplier survives in
amplitude while the support join still has only the values $0_L,1_L$.

The same formulas can be used without discarding the original mass.
For any supported nonzero scalar $\eta$, the diagonal gauge
$\operatorname{diag}(I,\eta I,\ldots,\eta^{d-1}I)$ gives
\[
 \eta\mathcal C_d(A)\cong\mathcal C_d(\eta^d A).
 \tag{FLW22}
\]
Indeed a forward step introduces one $\eta$ and the full return
introduces $\eta^d$. For $\eta=p^{k/2}$ this is precisely the
corner $p^{kd/2}A$, not $p^{k/2}A$. The gauge is not declared
unitary: its pulled-back block metric has the factors
$p^{kj}G_N$. Its power trace is exactly $p^{kn/2}$ times FLW13.

The signed odd family has an additional exact composition statement;
it must not be confused with Verschiebung at one fixed corner.
For an invertible original $F$ put
$\mathcal W_d(F)=V_d(\chi_d F^{\chi_d})$ for odd $d$.
For odd $c,d$ there is an explicit full-support isomorphism
\[
 \mathcal W_c(\mathcal W_d(F))\cong\mathcal W_{cd}(F).
 \tag{FLW36}
\]
Here the formula means an operator construction on the corresponding
direct sums of projectives; each newly introduced slot remains
independent. If $\chi_c=1$, flatten the two cyclic sums as before,
using $\chi_{cd}=\chi_d$.
If $\chi_c=-1$, write $A=\chi_dF^{\chi_d}$ and first compare
$-\mathcal C_d(A)^{-1}$ with $\mathcal C_d(-A^{-1})$.
The intertwiner from the latter to the former sends slot $j$ to
$-j\bmod d$ with block
$(-1)^j(A^{\lfloor-j/d\rfloor})^\uparrow$.
FLW12 with step $-1$ proves the inverse-cycle part. The successive
signs supply a factor $-1$ at each step, including the return since
$d$ is odd. Thus the map and its inverse are proved directly.
Its returning corner is
$-A^{-1}=-\chi_dF^{-\chi_d}=\chi_{cd}F^{\chi_{cd}}$.
Apply this same intertwiner in each of the $c$ outer slots and
then flatten the cycles; the result is FLW36.
In this inverse case its metric contains $G_N$ in slot zero and
$(A^{-1})^*G_NA^{-1}$ in each other slot, followed by the stated
permutation; the signs have absolute value one. It is not silently
made an isometry.

For general positive $n$, the exact Frobenius result for the signed
family is instead
\[
 F_n\mathcal W_d(F)
 \cong\bigoplus_{t=0}^{g-1}
 V_{d/g}(\chi_d^{n/g}F^{\chi_d n/g}),\qquad g=\gcd(d,n).
 \tag{FLW37}
\]
This is FLW18 with its literal corner substituted. For $(d,n)=1$
and odd $n$ it agrees with $\mathcal W_d(F^n)$.
For coprime even $n$ and $\chi_d=-1$, its corner is $F^{-n}$ whereas
the latter proposed replacement has corner $-F^{-n}$.
Their exact amplitude difference at that corner is $2F^{-n}$,
an invertible map, and both have the same cyclic support.
On the identified $d$ slots the signed difference of the two
endomorphisms has this last-to-first amplitude, while every
cyclic edge still retains its supported-zero support when its
other amplitudes cancel. Thus FLW37, rather than an unqualified
commutation assertion for the signed family, records the full
parity defect.

\subsection{Global synchronization and its exact defect object}

The single synchronized carrier of amplitude dimension $dq$ is
$P_{dq}$, whose free idempotent has all-top support in all $dq$
coordinates. It is not FLW5. The exact comparison consists of
\[
\begin{split}
 j:P_{dq}\longrightarrow P_{q,d},\quad
           (v,\lambda)&\longmapsto((v_0,\lambda),\ldots,(v_{d-1},\lambda)),\\
 r:P_{q,d}\longrightarrow P_{dq},\quad
           ((v_j,\lambda_j)_j)&\longmapsto
                         ((v_j)_j,\bigvee_j\lambda_j),\qquad rj=I.
\end{split}
\tag{FLW23}
\]
Substitution in FLW1 proves $S$-linearity. Their other composite
$jr$ synchronizes all slot labels to their join; it generally is
not the identity. The complete fibre over $(v,\lambda)$ is
\[
 \{(\lambda_j):\bigvee_j\lambda_j=\lambda,
                       \ v_j\ne0\Rightarrow\lambda_j=1_L\}.
 \tag{FLW24}
\]
Every vector amplitude is retained, and the exact identifications
among labels are displayed rather than removed silently.
The complex $C_d(A)$ has a lift $C_d(A)^\uparrow$ on $P_{dq}$,
with all-top support in every block. Because FLW11 permutes the
independent labels, the maps $j,r$ intertwine these two operators.
Its intrinsic trace is, at every integer $n$,
\[
 \operatorname{Tr}_{P_{dq}}((C_d(A)^n)^\uparrow)
                   =(\operatorname{tr}_{\mathbb C}C_d(A)^n,1_L).
 \tag{FLW25}
\]
At $d\nmid n$ FLW25 is $e$ whereas FLW13 is $\tau$.
For $d\mid n$ they agree exactly. Thus replacing independent cyclic
slots by the globally synchronized lift loses a proved part of the
semiring trace even though the entire original complex operator
and all its metric entries are the same.

This difference defines an actual comparison object: the split
surjection $r$, its section $j$, its idempotent $jr$, the full
fibres FLW24, and the two intertwined automorphisms.
For $d>1$ its label relation is nontrivial. For example choose any
slot and place support $1_L$ only there, with every amplitude zero;
all these distinct vectors map to $(0,1_L)$.
There is no assertion that an additive difference of semimodules
or a direct-sum complement of $j$ exists.

\subsection{What group completion kills, and what it does not kill}

The additive group completion of the coefficient semiring $S$ is
exactly $\mathbb C$. Indeed $z_\lambda+z_\lambda=z_\lambda$ forces
every $[z_\lambda]$ to be zero in the group. Every other element
has the form $(c,1_L)$, and
$(a,1_L)+(b,1_L)=(a+b,1_L)$.
The class of $(0,1_L)$ is zero, so these classes define the inverse
of the homomorphism induced by $\rho$ to the additive group
$\mathbb C$. Multiplication is compatible by FLW1. In particular
\[
 S^{\rm gp}\cong\mathbb C,\qquad [e]=[\tau]=0.
 \tag{FLW26}
\]
This proves a statement about scalar coefficients, not a vanishing
theorem for projective objects or their endomorphisms.

Here is a direct invariant detecting the distinction. Choose a
bounded lattice homomorphism $\beta:L\to\mathbb B=\{0,1\}$.
Such a homomorphism exists for a nontrivial bounded distributive
lattice: choose a maximal proper ideal by Zorn's lemma. It is prime,
because if $x,y$ are outside it, maximality gives $i\vee x=1$
and $j\vee y=1$ for ideal elements $i,j$; if $x\wedge y$ were in
the ideal, distributing these equalities would put $1$ in it.
The complementary indicator is the required bounded homomorphism.
This chosen probe does not replace $L$ in FLW1--FLW25.

Base change along $\beta s:S\to\mathbb B$ gives
\[
 (P_{q,d},\mathcal C_d(A))_{\mathbb B}
       =(\mathbb B^d,\text{cyclic permutation of the }d\text{ coordinates}),
 \qquad
 (P_{dq},C_d(A)^\uparrow)_{\mathbb B}=(\mathbb B,I).
 \tag{FLW27}
\]
Every returning block of $A^\uparrow$ has support all top, so it
acts as the identity on the single synchronized Boolean coordinate;
its nonzero amplitude, sign, and phase do not alter this map.
For the actual source, both objects are automorphisms, so their
stable images are their full underlying finite sets.

For any endomorphism $T$ of a finitely generated projective
$S$-semimodule $P$, its Boolean base change $P_{\mathbb B}$ is
finite: it is an image of $\mathbb B^a$ for a finite $a$.
Let $T_{\mathbb B}^{\infty}(P_{\mathbb B})$ be the eventual image
of its descending finite image sequence and define
\[
 b(P,T)=\log_2\big|T_{\mathbb B}^{\infty}(P_{\mathbb B})\big|.
 \tag{FLW28}
\]
It is invariant under isomorphism and additive on direct sums,
because direct sums of finite semimodules are their Cartesian
products and iteration acts componentwise. Null endomorphisms
have singleton eventual image and value zero. Thus FLW28 defines
a group homomorphism from the group generated by projective
endomorphism isomorphism classes with the direct-sum relations,
even after quotienting by all null-endomorphism classes. It also
respects the image-stabilization relation whenever that image is
a projective in this category. Write its factorization as
$T=ba$, with $a:P\to T(P)$ and $b:T(P)\to P$; the restricted
endomorphism on $T(P)$ is $ab$. After base change the same two
factorizations hold, even without assuming that the changed $b$
is injective. For two maps $a,b$ between finite sets, $a$ takes
the eventual image of $ba$ to that of $ab$ and intertwines their
permutations there. It is injective on that image because
$ax=ay$ implies $(ba)x=(ba)y$ and $ba$ is bijective there;
the same reasoning makes $b$ injective on the other eventual
image. The two finite eventual images therefore have equal
cardinality, proving the claimed relation.
For the two objects in FLW27 the values are respectively $d$ and $1$.
Their difference therefore has nonzero invariant $d-1$ for $d>1$.
No relation from a different, unspecified exact category is assumed.

Even a projective having zero complex amplitude can survive this
group. Namely $eS=\operatorname{im}(e:S\to S)$ is projective
because $e^2=e$. Its amplitude base change is zero, its full support
module is $L$, and its Boolean base change is $\mathbb B$.
Consequently
\[
 \operatorname{Tr}_{eS}(I)=e,\qquad
 [e]=0\text{ in }S^{\rm gp},\qquad b(eS,I)=1.
 \tag{FLW29}
\]
These explicit maps prove that scalar group completion kills its
trace coefficient while the projective endomorphism class remains
nonzero in the stated additive quotient.

\subsection{The exact pointed-cycle receiver and corrected multiplicities}

The support module before choosing $\beta$ is $L^d$ with the literal
cyclic permutation of its slots. After the probe it is $\mathbb B^d$,
whose nonzero join-irreducible elements are precisely the $d$ unit
coordinate vectors: a subset with at least two elements is a join
of two proper subsets, whereas a singleton has no such expression.
An automorphism preserves this property and therefore permutes the
same $d$ elements. Adding a fixed base point gives the finite pointed
cycle $C(d)$ in the notation of Connes--Consani.
For free Boolean modules, this operation respects direct sums and
tensor products: singleton coordinates of $\mathbb B^{a+b}$ are
the disjoint union of those of its factors, and the free tensor
$\mathbb B^a\otimes\mathbb B^b=\mathbb B^{ab}$ has coordinate
pairs. These identifications intertwine the permutations.
Thus it gives an exact additive and multiplicative comparison on
the subcategory of the support cycles constructed here.

The original pointed-cycle ghost and the scalar lattice trace now
have their precise relation:
\[
 \operatorname{gh}_n C(d)=d\,\mathbf1_{d\mid n},\qquad
 \operatorname{tr}_{L}(P_d^n)=
       \begin{cases}1_L,&d\mid n,\\0_L,&d\nmid n.
       \end{cases}
 \tag{FLW30}
\]
The first counts the $d$ fixed nonbasepoint slots; the second is
the join of their diagonal indicators. It is obtained from the
nonnegative integer ghost by the homomorphism
$\mathbb N\to L$ sending zero to $0_L$ and every positive integer
to $1_L$. The actual cycle object retains the multiplicity that
this homomorphism loses. For comparison the full fixed set of
$P_d^n$ on $\mathbb B^d$ has cardinality
$2^{\gcd(d,n)}$: a fixed subset is a union of the
$\gcd(d,n)$ permutation orbits. Thus counting all subset states
would give a different invariant from the pointed atom trace;
the explicit atom map specifies the source invariant being used.

The root-of-unity receiver of the authors' theorem is exactly
\[
 \Theta_d=\sum_{d\alpha=0}[\alpha]\in\mathbb Z[\mathbb Q/\mathbb Z],
 \qquad
 \widehat\Theta_d(n)=\sum_{j=0}^{d-1}e^{2\pi i nj/d}
                   =d\mathbf1_{d\mid n}.
 \tag{FLW31}
\]
The finite geometric sum proves the equality. Hence the full-support
comparison defect FLW23 has pointed-cycle image
$\Theta_d-\Theta_1$, with ghost $d\mathbf1_{d\mid n}-1$.
It is nonzero for $d>1$, for example at $n=1$.
This proves its connection with the authors' universal invariant
without asserting that the original semiring scalar trace itself
is that group-ring invariant.

There are two source formulas that require explicit correction in
the cited original version. If $m(k)$ denotes the number of cycles
of length $k$, direct counting gives
\[
 g_n=\sum_{k\mid n}k\,m(k),\qquad
 m(k)=\frac1k\sum_{d\mid k}\mu(k/d)g_d.
 \tag{FLW32}
\]
For the inverse formula, substitute the first sum. The coefficient
of $j\,m(j)$ is $\sum_{j\mid d\mid k}\mu(k/d)$, which equals one
when $j=k$ and zero otherwise. This divisor identity follows by
expanding $\prod_{p\mid(k/j)}(1-1)$, with value one for $k/j=1$.
The source at lines 510--512 omits both the cycle length in its
preceding count and the factor $1/k$ in its inverse. A single
two-cycle has $g_1=0,g_2=2$ and $m(2)=1$, which checks the factor
directly. The primitive-root formula at line 516 has a mismatched
dummy index; its correctly indexed form is
$\sum_{d\mid n}\mu(n/d)\Theta_d$, following from partitioning
roots by exact order. These corrections do not alter the cycle
invariant FLW31.

Even the corrected signed M\"obius coefficient over the full
semiring has a retained supported-zero tail. For the operator
$\mathcal C_d(I)$ on $P_{q,d}$, let $t_n$ denote its FLW13 trace,
and form, for $k\ge1$, the literal finite sum
\[
 M_k^S=(1/k,1_L)\sum_{\substack{l\mid k\\\mu(k/l)\ne0}}
                  (\mu(k/l),1_L)t_l
 =\begin{cases}
 \tau,&d\nmid k,\\
 (q,1_L),&k=d,\\
 e,&d\mid k,\ k\ne d.
 \end{cases}
 \tag{FLW33}
\]
Its amplitude is FLW32 for $q$ complex copies of the $d$-cycle.
If $d\mid k$, the summand $l=k$ with M\"obius value one has
top support; the support join therefore remains top even when all
complex terms cancel. If $d\nmid k$, none of its divisors returns
and every summand is $\tau$. This proves all cases of FLW33 and
exhibits precisely why subtraction after scalar group completion
loses more information than the original semiring calculation.

\subsection{Adams operations and the actual domain of Newton recovery}

For the original amplitude corner $A=\chi F_p^\chi$, retain its
labelled eigenvalues $\alpha_{ab}=\chi\omega_{ab}^\chi$ and set
$\lambda_A(t)=\det(I+tA)=\prod_{a,b}(1+\alpha_{ab}t)$.
Direct formal differentiation of this finite product proves
\[
 t\lambda_A'(t)/\lambda_A(t)
   =\sum_{n\ge1}(-1)^{n-1}\operatorname{tr}(A^n)t^n,
 \qquad
 j e_j=\sum_{n=1}^j(-1)^{n-1}e_{j-n}\operatorname{tr}(A^n),
 \tag{FLW34}
\]
where $\lambda_A(t)=\sum_{j=0}^q e_jt^j$ and $e_0=1$.
For the first identity expand
\[
\frac{\alpha t}{1+\alpha t}=\sum_{n\ge1}(-1)^{n-1}\alpha^nt^n;
\]
for the second multiply by $\lambda_A(t)$ and compare coefficients.
The sign at \texttt{lognewton}, source lines 644--647, is printed
as $(-1)^n$; the coefficient of $t$ already requires the displayed
$(-1)^{n-1}$. This agrees with the source's preceding explicit
$\psi_1=\lambda_1$ and $\psi_2=\lambda_1^2-2\lambda_2$.

For the cyclic amplitude receiver,
\[
 \det(I-tC_d(A))=\det(I-t^dA)
            =\prod_{a,b}(1-\chi\omega_{ab}^{\chi}t^d),
 \qquad
 \operatorname{tr}(A^m)=d^{-1}\rho
                       (\operatorname{Tr}\mathcal C_d(A)^{dm}).
 \tag{FLW35}
\]
The determinant follows by restricting to each original Lagrange
label, whose cyclic matrix has polynomial $z^d-\alpha_{ab}$.
Those labelled invariant complex subspaces sum to the full amplitude
space, so all multiplicities are retained. The second equality is
FLW13. FLW34 then recovers every coefficient of the original
$\det(I-tA)$ from its first $q$ returning traces. For the weighted
receiver the factor is $\chi p^{kd/2}\omega_{ab}^\chi$ and the
returning traces include $p^{kdm/2}$ exactly as in FLW14--FLW15.

FLW34--FLW35 are statements in the complex amplitude ring or its
formal series ring. They are not assumed determinant or exterior
power identities for the full semiring projectives. Over $S$,
$(1,1_L)+(-1,1_L)=e\ne\tau$; an alternating complex cancellation
does not give the semiring additive zero. The full coefficient
object retained in parallel is FLW15 together with its projective
slot action FLW11 and the comparison FLW23--FLW24. Its exact images
under both coefficient maps and the surviving cycle invariant have
all been proved above. In particular neither the Newton amplitude
recovery nor scalar group completion deletes the constructed
projective support defect.
